\section{Introduction}
\label{sec:intro}

When a modern LLM solves a benchmark item, we usually cannot tell whether it generalized from underlying competence or simply retrieved a near-copy of the answer from its pretraining corpus. \citet{wang2025generalization} show this matters quantitatively: for factual tasks, accuracy is largely explained by n-gram co-occurrence frequencies in pretraining data, not by deeper competence. Distinguishing the two currently requires expensive contamination audits over multi-terabyte corpora~\citep{lee2021deduplicating,wu2024antileak} or carefully constructed post-cutoff benchmarks. Both routes are slow, partial, and must be redone for every new model.

\para{The Talkie shortcut.} \citet{talkielm2026} recently released \talkienineteen, a 13B-parameter decoder-only LLM trained on 260B tokens of strictly pre-1931 English (books, periodicals, scientific journals, patents, case law) together with an architecture- and FLOPs-matched modern twin \talkieweb trained on FineWeb. The 1930 cutoff is a structural decontamination guarantee: every post-1930 concept --- Python, the iPhone, CRISPR, the Hiroshima bombing --- has co-occurrence probability \emph{zero} in the corpus by construction. Any non-trivial \talkienineteen performance on a post-1930 task must therefore be generalization.

\para{The meta-hypothesis.} Holtzman's hypothesis (submitted to the IdeaHub one week before our experiments) sharpens this further: because \talkienineteen's corpus covers human experience much more sparsely than modern web data, the generalization signal-to-noise ratio is \emph{qualitatively} better. The claim is not "fewer items needed for significance" but "items become decisively interpretable rather than ambiguous." This is the meta-claim we test --- a methodological claim about evaluation infrastructure, not a model-quality claim.

\para{Our approach.} We run three matched probes on the triple \{\talkienineteen, \talkieweb, \gptfour\}: a 50-item post-1930 probe battery (factual completion plus in-context Python/JavaScript items), a date-stratified \mmlu subsample, and \humaneval. On top of these we layer a Bayesian attribution-bits framework~(\secref{sec:method}) that quantifies how many bits of evidence each correct answer contributes to "the model generalized." For \talkienineteen on post-1930 items the prior probability of memorization is approximately zero by construction, so the framework yields large positive bits per correct answer; for the modern baselines on the same items the prior is unknown and bits collapse to $\le 0$.

\para{Quantitative preview.} \talkienineteen attains 3\% post-1930 fact recall versus 73\% (\talkieweb) and 80\% (\gptfour), with a paired per-token surprisal gap of +6.84 \bpt over its modern twin (paired permutation $p < 0.001$, $n=30$). On a single Python in-context example, \talkienineteen correctly continues \texttt{def triple(x): return } $\to$ \texttt{x * 3} and \texttt{cubes = [x**3 for x in range(5)] \# cubes is } $\to$ \texttt{[0, 1, 8, 27, 64]}, despite zero Python exposure in training. Each such success is, by construction, uncontaminated evidence of generalization worth $\approx 34$ bits per item under a worst-case attribution prior; the same accuracy on the modern baselines yields $\le 0$ bits without a corpus audit.

\para{Contributions.}
\begin{itemize}[leftmargin=*,itemsep=0pt,topsep=0pt]
    \item We operationalize Holtzman's meta-hypothesis as a Bayesian attribution-bits framework that converts the structural pre-1931 cutoff into a quantitative argument: each correct \talkienineteen answer on a post-1930 item is worth $\approx 34$ bits, against $\le 0$ bits for modern baselines under a worst-case prior.
    \item We construct a 50-item post-1930 probe battery (30 facts, 9 pre-1930 controls, 11 in-context Python/JS items) that exposes the cutoff cleanly: the paired surprisal gap between \talkienineteen and its matched modern twin is +6.84 \bpt ($p < 0.001$, $n=30$).
    \item We document positive evidence of uncontaminated in-context generalization: \talkienineteen succeeds on 4 of 11 ICL items despite zero exposure to Python/JavaScript, an unusually clean positive-evidence example for generalization in current LLMs.
    \item We show where the cutoff design does \emph{not} help (date-stratified \mmlu is too noisy at single-question granularity; bare-base 13B models fail \humaneval for alignment reasons unrelated to the cutoff) and explain why, isolating where the methodology lands its load.
\end{itemize}

The rest of the paper is organized as follows. \Secref{sec:related} positions our work against generalization-vs-memorization analyses and contamination-controlled benchmarks. \Secref{sec:method} describes the models, probes, scoring, and the attribution-bits framework. \Secref{sec:results} reports the three probes and the meta-analysis. \Secref{sec:discussion} discusses limitations and threats to validity. \Secref{sec:conclusion} concludes.
